\begin{lemma}
Let \(0<\eta<\varepsilon\).  For all sufficiently large \(n\), the rounded
integer scale
\[
K_{1+\eta}(n)=
\noderef{TwoBiteNaturalIndependenceScale}(1+\eta,n)
\]
is still strictly below the larger real target
\[
\noderef{TwoBiteIndependenceScale}(1+\varepsilon,n)
=(1+\varepsilon)\sqrt{n\log n}.
\]
\end{lemma}
\begin{proof}
Set
\[
s(n)=\sqrt{n\log n},\qquad
x_\eta(n)=(1+\eta)s(n),\qquad
x_\varepsilon(n)=(1+\varepsilon)s(n).
\]
From \(0<\eta\) and \(\eta<\varepsilon\), we have
\[
0<\varepsilon-\eta
\qquad\text{and}\qquad
0<1+\eta.
\]
The gap between the two real scales is
\[
x_\varepsilon(n)-x_\eta(n)
=(\varepsilon-\eta)s(n).
\]
The factor \(\varepsilon-\eta\) is fixed and positive, while
\(s(n)=\sqrt{n\log n}\to\infty\).  Choose \(n_0\) so large that for every
\(n\ge n_0\):
\[
\log n\ge0, \tag{1}
\]
\[
0\le \noderef{TwoBiteReducedVertexCount}(n)
\quad\text{and}\quad
0\le \noderef{TwoBiteIndependenceScale}(1+\eta,n), \tag{2}
\]
and
\[
1<(\varepsilon-\eta)s(n). \tag{3}
\]
The first inequality ensures \(s(n)\) is the square root of a nonnegative
quantity.  Together with \(0<1+\eta\), it gives the second nonnegativity
condition in (2); the reduced-count nonnegativity follows from
\(\noderef{TwoBiteReducedVertexCount}(n)=n/\log^2 n\) for the same large \(n\).

Now fix \(n\ge n_0\).  Apply \(\noderef{NaturalParameterApproximation}\) with
\(\kappa=1+\eta\), using the two nonnegativity hypotheses in (2).  Its fourth
conjunct gives the rounded-ceiling upper bound
\[
(K_{1+\eta}(n):\mathbb R)<x_\eta(n)+1.
\]
Indeed, after unfolding \(\noderef{TwoBiteIndependenceScale}\), the right side
is
\[
\noderef{TwoBiteIndependenceScale}(1+\eta,n)+1
=(1+\eta)s(n)+1=x_\eta(n)+1.
\]
By (3),
\[
x_\eta(n)+1<x_\eta(n)+(\varepsilon-\eta)s(n).
\]
Unfolding the two real scales and collecting the coefficient of \(s(n)\),
\[
x_\eta(n)+(\varepsilon-\eta)s(n)
=(1+\eta)s(n)+(\varepsilon-\eta)s(n)
=(1+\varepsilon)s(n)
=x_\varepsilon(n).
\]
Therefore
\[
(K_{1+\eta}(n):\mathbb R)
<x_\eta(n)+1
<x_\varepsilon(n)
=\noderef{TwoBiteIndependenceScale}(1+\varepsilon,n).
\]
This proves the asserted strict comparison for every \(n\ge n_0\).
\end{proof}
